\section{Conclusion}
\label{sec:conclusion}

The empirical findings of this study affirm one half of the central thesis that a
regime-sensitive framework for signal integration and portfolio construction delivers an
improved risk-return dynamic compared to static approaches, and refute the other. Superior
capital retention during market downturns, a lower volatility signature in bullish phases,
and a rolling Sharpe ratio positive in the majority of windows demonstrate the efficacy of
one of the model's key components---its ability to adjust total risk to evolving market
conditions---together with an implementation strategy that accounts for transaction costs.
They do not support its probabilistic fusion of diverse signal types. This paper set out
to close the step between a predictive signal and an executable portfolio, and to
determine which parts of a regime-conditional framework are responsible for its behaviour.
Both objectives are met, though not in the direction originally anticipated.

Evaluated on ten liquid multi-asset instruments over eighteen years under a walk-forward
protocol with transaction costs, the framework raises the Sharpe ratio from 0.628 to
0.878 and reduces maximum drawdown from $-36.1$ to $-19.0$ per cent against an
equal-weighted benchmark, while earning a lower compound return. The advantage holds in
all four stress episodes, in 18 of 19 calendar years on drawdown, in all eighteen
specifications examined, and in a held-out final block that postdates every specification
decision---but it is consistent rather than statistically significant, its bootstrap
interval including zero. Ablation locates its source: at identical average exposure,
timing risk by the identified market state improves both measures with no difference in
mean return, and the form of the risk budget proves more consequential than its
calibration, since the absolute targets of the original specification bind in one
rebalance out of 224 and leave the layer inert.

The signal fusion layer, which gave the framework its name, does not contribute. Removing
it improves the Sharpe ratio, the information coefficients of both composites are
significantly negative over the full sample, and the regime-conditional weighting schedule
has no measured basis, since the calm-to-stress contrast it presumes is not significant and
its point estimates run the other way. Nor is the negative relationship stable enough to
invert, since it disappears in the holdout. This null is reported in full, with the
coefficients that produce it, because a component-level negative result is more useful than
an aggregate number that conceals which parts of an architecture are doing the work.

The study's contribution is therefore a specification and a method of testing it: a
two-step separation of composition from total risk that makes each layer independently
falsifiable, an ablation design that isolates risk timing from risk level by matching
average exposure, and a documented account of two failure modes---an inert absolute risk
budget, and a signal-to-allocation conversion in which mismatched units rather than
information determine the solution. The framework is useful within stated bounds: it
requires a universe with genuine dispersion in asset volatility, trades return for
drawdown at a measured rate, and should be evaluated against a volatility-matched rather
than an unlevered benchmark. Outside those bounds, as the mega-capitalisation results
show, it does not add value. Performance of this kind across varied and often challenging
market scenarios highlights the necessity of addressing the theoretical disconnect between
signal generation and allocation, and the study offers both a methodological foundation
for a new generation of adaptive investment models and an empirical account of what such
models can and cannot be shown to achieve---provided each layer of the resulting
architecture is tested separately rather than credited with the performance of the whole.
